\section{Quasiparticle corrections}\label{app:E.3}

For an electronic threshold the correction of the Kohn--Sham gap, more than half an eV in silicon (\cref{fig:3-si-qp}), is the largest a calculation makes. The $GW$ self-energy $\Sigma = \ii\GF W^{\mathrm{scr}}$, with $W^{\mathrm{scr}} = \varepsilon^{-1}v_{\mathrm{C}}$, and the electron--electron lifetime from $\Imag\Sigma$ are defined in \cref{sec:4.8}. The quasiparticle corrections run from the scissor shift to a one-shot $G_0W_0$ calculation and its self-consistent successors, and each changes the energies, the lifetimes and the matrix elements of a dark-matter rate.

\subsection{Computed quantities}

The quantities are the quasiparticle energies $E_{\mathrm{qp}}(n\vect{k})$, the positions of the poles of the one-particle Green's function $\GF$ (\cref{sec:4.7}), with their residues $z_{\mathrm{qp}}$ and their widths. In practice they are computed as corrections to the Kohn--Sham energies, and the Kohn--Sham states are kept. A quasiparticle calculation changes the inputs of every later step (the threshold $E_{\mathrm{min}}(\vect{q})$, the energy denominators of a response, the screening) but not the approximation level of that step.

\subsection{The scissor shift in practice}

The accuracy of the scissor shift \cref{eq:3.4-scissor} in silicon, about \SI{0.1}{eV} across the bands, is partly fortuitous: the energy dependence of $\Sigma$ and its nonlocality would each make the quasiparticle correction strongly state-dependent, and in silicon the two almost cancel~\cite[Sect.~V\,B]{aryasetiawan1998gw}. The scissor shift raises the $f$-sum of the response by $\Delta_{\mathrm{scis}}\sum_c|\rho_{cv}|^2$, and this excess is the extra term \cref{eq:E.2-fsumNL} generated by a nonlocal operator. A scissor-corrected response therefore violates the $f$-sum rule unless the transition strengths are rescaled with the energies, and a response used in an $f$-sum check must be built consistently (\cref{app:E.6}).

The shift also changes the screening. In the independent-particle response each transition contributes $|\rho_{cv}|^2$ divided by its energy, so with $\rho_{cv}$ fixed and the energies raised, $\varepsilon_\infty - 1$ decreases. The PBE value of silicon, 13.04 against the measured 11.7, moves in the direction of the measurement. Whether the scissored value lands above or below 11.7 depends on the rest of the calculation, local fields included, so a scissored screening is not guaranteed to be closer to the measured one.

Shifting a static-lattice calculation to the measured gap and then adding the electron--phonon renormalization of \cref{sec:13.7,app:E.8} counts the renormalization twice; in silicon the zero-point part is about \SI{60}{meV} (\cref{we:E-zpr}), small next to the gap but not against the window widths of \cref{chap:13}. An optical absorption edge lies below the fundamental gap by the exciton binding energy (\cref{sec:4.9}), and a scissor to the optical edge puts the free-pair continuum too low.

\subsection{One-shot \texorpdfstring{$G_0W_0$}{G0W0}}

The standard method for semiconductors is the one-shot $G_0W_0$ method of Hybertsen and Louie~\cite{hybertsen1986electron}: $\GF$ and $W^{\mathrm{scr}}$ are built once from the Kohn--Sham states, and the energies are corrected to first order.

The input is a Kohn--Sham calculation, extended to many empty bands, of the order of a hundred per atom (the count is made below). The screening is built from the static independent-particle response $R^0_{\vect{G}\vect{G}'}(\vect{q}, \omega = 0)$ of \cref{sec:4.6}. $R^0$ is a sum over pairs of occupied and empty states, and each term is a product of the plane-wave matrix elements
\begin{equation}
  \rho_{mn}(\vect{k}, \vect{q} + \vect{G}) = \langle\psi_{m,\vect{k}+\vect{q}}|\ee^{\ii(\vect{q}+\vect{G})\cdot\vect{r}}|\psi_{n\vect{k}}\rangle ,
  \label{eq:E.3-M}
\end{equation}
divided by the energy difference of the pair~\cite[Eqs.~(7) and (8)]{deslippe2012berkeleygw}. These matrix elements are the density vertices of \cref{sec:3.2} at the unreduced momentum $\vect{q} + \vect{G}$ (in the periodic gauge, $\psi_{m,\vect{k}+\vect{q}+\vect{G}} = \psi_{m,\vect{k}+\vect{q}}$). The density vertices that enter a dark-matter rate also enter the screening of the electron's own interaction, and the large-$q$ caveats of \cref{app:E.2} apply to the rate and to the screening alike. From $R^0$ follow the dielectric matrix of the random-phase approximation, $\varepsilon_{\vect{G}\vect{G}'} = \delta_{\vect{G}\vect{G}'} - v_{\mathrm{C}}(\vect{q} + \vect{G})R^0_{\vect{G}\vect{G}'}$ (\cref{sec:4.6}), its inverse, and the screened interaction $W^{\mathrm{scr}}_{\vect{G}\vect{G}'}(\vect{q}) = (\varepsilon^{-1})_{\vect{G}\vect{G}'}(\vect{q})\,v_{\mathrm{C}}(\vect{q} + \vect{G}')$~\cite[Eq.~(12)]{deslippe2012berkeleygw}.

The frequency dependence of $W^{\mathrm{scr}}$ comes from a model or from a frequency grid. In the generalized plasmon-pole model of Hybertsen and Louie, each element of $\Imag(\varepsilon^{-1})_{\vect{G}\vect{G}'}(\vect{q}, \omega)$ at $\omega > 0$ is a single delta function at a frequency $\tilde\omega_{\vect{G}\vect{G}'}(\vect{q})$, and the real part follows by the Kramers--Kronig relation. The two parameters of each element, the weight and the frequency of its pole, are fixed by two exact constraints: the static value $(\varepsilon^{-1})_{\vect{G}\vect{G}'}(\vect{q}, 0)$ just computed, and the $f$-sum rule for $\varepsilon^{-1}$, $\int_0^\infty\dd\omega\,\omega\,\Imag(\varepsilon^{-1})_{\vect{G}\vect{G}'} = -(\pi/2)\,\omega_p^2\,[(\vect{q} + \vect{G})\cdot(\vect{q} + \vect{G}')/|\vect{q} + \vect{G}|^2]\,\nd(\vect{G} - \vect{G}')/\nd(0)$, with $\nd(\vect{G})$ the Fourier components of the density~\cite[Eqs.~(83)--(86)]{aryasetiawan1998gw}. No parameter is adjusted. The model works because in sp semiconductors most of the spectral weight of $\Imag\varepsilon^{-1}$ is in the plasmon, near the free-electron plasma frequency of the valence electrons, which for the eight valence electrons in the \SI{40.05}{\angstrom^3} cell of silicon is $\omega_p = \SI{16.6}{eV}$. The alternative is to compute $\varepsilon^{-1}$ at many frequencies, on the real axis or on the imaginary axis followed by an analytic continuation. The continuation is usually a rational (Pad\'e) fit, and \cref{app:C.5} explains why rational fits are unstable and must be checked against sum rules.

The self-energy is a sum over bands, $\vect{q}$ and $\vect{G}$ of products of the matrix elements \cref{eq:E.3-M}, weighted by $v_{\mathrm{C}}$ and by the frequency structure of $\varepsilon^{-1}$. It is organized as a screened-exchange part, a sum over occupied states only, and a Coulomb-hole part, a sum over all bands, occupied and empty~\cite[Eqs.~(19)--(22)]{deslippe2012berkeleygw}. Codes evaluate two parts separately: the bare exchange $\Sigma_{\mathrm{x}}$, which is the Fock exchange of \cref{sec:3.4} built from the occupied Kohn--Sham states, and the correlation part $\Sigma_{\mathrm{c}}$ that $W^{\mathrm{scr}} - v_{\mathrm{C}}$ adds to $\Sigma_{\mathrm{x}}$.

The quasiparticle equation, which corrects the energies, replaces $v_{\mathrm{xc}}$ in the Kohn--Sham equation by $\Sigma(E)$, evaluated at the quasiparticle energy itself. Keeping only the diagonal matrix elements in the Kohn--Sham basis assumes that the Kohn--Sham states are close to the quasiparticle states. In this diagonal approximation, expanding $\Sigma(E)$ to first order about $\varepsilon_{n\vect{k}}$ gives~\cite[Eqs.~(57) and (58)]{aryasetiawan1998gw}
\begin{equation}
  E_{\mathrm{qp}}(n\vect{k}) = \varepsilon_{n\vect{k}} + z_{\mathrm{qp}}\,\Real\langle\psi_{n\vect{k}}|\Sigma(\varepsilon_{n\vect{k}}) - v_{\mathrm{xc}}|\psi_{n\vect{k}}\rangle, \qquad
  z_{\mathrm{qp}} = \biggl[1 - \frac{\partial\Real\langle\psi_{n\vect{k}}|\Sigma(\omega)|\psi_{n\vect{k}}\rangle}{\partial\omega}\biggr]^{-1}_{\omega = \varepsilon_{n\vect{k}}} .
  \label{eq:E.3-qp}
\end{equation}
The residue $z_{\mathrm{qp}}$ is 0.7 to 0.8 in semiconductors~\cite[Sect.~V\,B]{aryasetiawan1998gw}. The slope of $\Real\Sigma$ is negative, so $z_{\mathrm{qp}} < 1$ reduces the correction: neglecting the energy dependence of $\Sigma$, as the static Coulomb-hole plus screened-exchange (COHSEX) approximation does, overestimates the gap correction~\cite[Sect.~V\,B]{aryasetiawan1998gw}. The off-diagonal elements, when kept, turn \cref{eq:E.3-qp} into a small matrix diagonalization in the Kohn--Sham basis~\cite[Eqs.~(17) and (18)]{deslippe2012berkeleygw}, which changes the states as well as the energies.

\begin{table}[t]
  \centering
  \caption{Minimum band gaps of silicon, germanium and GaAs in LDA, in one-shot $GW$ calculations started from LDA, and from experiment, in eV, as collected in Ref.~\cite[Table~I]{aryasetiawan1998gw}. The $GW$ column gives the range of the independent calculations of Refs.~\cite{hybertsen1986electron,godby1988self,rohlfing1993quasiparticle,zhu1991quasiparticle}. The LDA gap of germanium is negative (the crystal is a metal); for GaAs the LDA gap is the one of \cref{tab:E-functionals}, \SI{0.43}{eV}. The last two columns are the $GW$ calculations of Shirley, Zhu and Louie~\cite{shirley1992core}, in which the core--valence exchange and correlation are treated with the local-density approximation (``core in LDA'') or with a core-polarization potential, as collected in Ref.~\cite[Table~II]{aryasetiawan1998gw}: the silicon gap from $\Gamma$ to the conduction minimum at $0.85\,$X, and the direct ($\Gamma$) and indirect (L) gaps of germanium, with spin--orbit coupling. The measured direct gap of germanium is \SI{0.89}{eV}~\cite[Table~II]{aryasetiawan1998gw}.}
  \label{tab:E-gw}
  \begin{tabular}{l ccc cc}
    \toprule
    & LDA & $GW$ & measured & $GW$, core in LDA & $GW$, core polarization \\
    \midrule
    Si & 0.52 & 1.24--1.29 & 1.17 & $0.85\,$X 1.29 & $0.85\,$X 1.16 \\
    Ge & $<0$ & 0.65--0.75 & 0.744 & $\Gamma$ 0.53, L 0.58 & $\Gamma$ 0.85, L 0.73 \\
    GaAs & -- & 1.22--1.58 & 1.52 & -- & -- \\
    \bottomrule
  \end{tabular}
\end{table}

\Cref{tab:E-gw} shows what these $G_0W_0$ calculations achieve for silicon, germanium and GaAs. In silicon they correct most of the LDA error and overshoot the measured gap by 0.07 to \SI{0.12}{eV}. In germanium and GaAs the independent calculations scatter by 0.1 and \SI{0.36}{eV}. The scatter measures the sensitivity of the method to its numerical and physical choices rather than its intrinsic accuracy. One of these choices is the treatment of the core. The calculation of Shirley, Zhu and Louie~\cite{shirley1992core} includes spin--orbit coupling and resolves the direct and indirect gaps of germanium. When that calculation treats the exchange and correlation of the valence electrons with the 3d and deeper shells of germanium within LDA, it puts the direct gap below the indirect one, the wrong order. Treating the core polarization explicitly restores the order, brings both gaps within \SI{0.04}{eV} of experiment, and lowers the silicon gap from 1.29 to \SI{1.16}{eV} (\cref{tab:E-gw}, last two columns). The core electrons must therefore be treated beyond LDA in the self-energy, either in the valence or through a core-polarization model, and not only in the density.

\subsection{Self-consistency and the starting point}

A one-shot calculation inherits its starting point. $G_0W_0$ on LDA states underestimates gaps systematically once core states are treated adequately, and it can fail qualitatively in transition-metal compounds~\cite{vanschilfgaarde2006quasiparticle}. Updating the energies in $\GF$ and $W^{\mathrm{scr}}$ until they are self-consistent (eigenvalue self-consistency) removes part of the dependence on the starting point. The quasiparticle self-consistent $GW$ method of van Schilfgaarde, Kotani and Faleev goes further and determines the best static, Hermitian one-particle Hamiltonian from the self-energy itself, so that the result no longer depends on the starting functional~\cite{vanschilfgaarde2006quasiparticle}. Its errors are small and systematic. It overestimates the gaps of sp semiconductors by about \SI{0.2}{eV} for transitions between $\Gamma$ points and \SI{0.1}{eV} for the others, and it slightly underestimates $\varepsilon_\infty$. The reason given is that the random-phase approximation of $W^{\mathrm{scr}}$ omits the electron--hole attraction, and without that attraction the screening is too weak and the gap too large~\cite{vanschilfgaarde2006quasiparticle}. A starting point that is metallic, LDA or PBE germanium, is the case where the choice matters most: the first $R^0$ then contains intraband transitions that an insulator does not have, and either a hybrid-functional start or a self-consistent scheme is used.

Current practice differs from the calculations of \cref{tab:E-gw} in three ways: the convergence of the basis, the treatment of the frequency integral, and the vertex~\cite[Sects.~IV\,C--F, V and~XI\,D]{golze2019gw}. The quasiparticle energies converge only as the inverse of the number of plane waves in the basis for the empty states and the response, so that the few hundred empty bands per atom of the early calculations can leave errors of tenths of an eV. Modern calculations extrapolate in $1/N_{\mathrm{PW}}$, or avoid the sum over empty states with a Sternheimer equation or a static remainder; converged in this way, $G_0W_0$ started from LDA gives \SI{1.10}{eV} for silicon~\cite[Sect.~IV\,C]{klimes2014predictive}, below the measured gap rather than above it. The frequency integral is done by contour deformation or on the imaginary axis rather than with a plasmon pole. And the missing electron--hole attraction in $W^{\mathrm{scr}}$ is put back as a vertex correction. The ladder diagrams of the Bethe--Salpeter equation (\cref{sec:4.9}), added to the polarizability of quasiparticle self-consistent $GW$, remove most of the overestimate of the gaps and the underestimate of $\varepsilon_\infty$ by that method. With the ladder diagrams, the silicon gap near X is \SI{1.08}{eV} against the measured \SI{1.17}{eV}~\cite[Table~1]{cunningham2023qsgw}.

\subsection{Lifetimes}

The imaginary part of the self-energy gives the width of the quasiparticle. Near the pole, the linearized quasiparticle energy of \cref{eq:E.3-qp} acquires an imaginary part $z_{\mathrm{qp}}\Imag\Sigma$, and the probability that the electron is still in the state decays as $\exp(-2z_{\mathrm{qp}}|\Imag\Sigma|\,t)$. The decay rate of the electron, its inverse lifetime, is therefore $2z_{\mathrm{qp}}|\Imag\Sigma(E_{\mathrm{qp}}(n\vect{k}))|$. By the cutting rule of \cref{sec:4.8}, $\Imag\Sigma$ in the $GW$ approximation is the golden-rule rate of the decay of the electron into a lower state with the emission of an excitation of the screened interaction: an electron--hole pair or a plasmon. An electron in the conduction band of an insulator can create an electron--hole pair only if it has the energy to cross the gap. Its energy above the conduction-band minimum must exceed $\Eg$, and the conservation of crystal momentum raises the threshold for pair creation above $\Eg$. Below that threshold the electron--electron width vanishes, and only phonon emission (\cref{sec:6.3}) relaxes the electron. Above the threshold, impact ionization opens, and its rate grows with the phase space of the final pairs. The phonon-emission rate and the impact-ionization rate, as functions of energy, are the input of the hot-carrier cascades that set the mean energy per electron--hole pair $\varepsilon_{eh}$ (\cref{app:E.13}).

The plasmon-pole model, with all the weight of $\Imag\varepsilon^{-1}$ in delta functions at the plasmon energies, gives $\Imag\Sigma = 0$ except at those energies, and it cannot produce lifetimes~\cite[Sect.~III\,C]{aryasetiawan1998gw}. A lifetime needs the full frequency dependence of $\varepsilon^{-1}$ near the energies of the created pairs; in silicon these energies start at the gap. A full-frequency calculation costs more than a plasmon-pole one by about the number of frequencies at which $\varepsilon^{-1}$ is computed.

\subsection{Matrix elements}

The diagonal approximation of \cref{eq:E.3-qp} keeps the Kohn--Sham states and changes the energies only, and its effect on matrix elements follows from the unchanged states and the new energies alone. The density vertex $\rho_{mn}(\vect{k}, \vect{q})$ is unchanged. The velocity changes, and its change can be computed exactly within the approximation. For eigenstates of a Hamiltonian $H$, the interband velocity is $\vect{v}_{mn} = \ii\langle m|[H, \vect{r}]|n\rangle = \ii(\varepsilon_m - \varepsilon_n)\langle m|\vect{r}|n\rangle$ for $m \neq n$. Between Bloch states at the same $\vect{k}$ the interband position matrix element is well defined, and this relation, applied to the Kohn--Sham Hamiltonian, gives the position matrix element as $\langle m|\vect{r}|n\rangle = -\ii\vect{v}^{\mathrm{KS}}_{mn}/(\varepsilon_m - \varepsilon_n)$. A quasiparticle Hamiltonian that is diagonal in the Kohn--Sham basis has the same $\langle m|\vect{r}|n\rangle$ and the energies $E_{\mathrm{qp}}$, so
\begin{equation}
  \vect{v}^{\mathrm{qp}}_{mn}(\vect{k}) = \vect{v}^{\mathrm{KS}}_{mn}(\vect{k})\,\frac{E_{\mathrm{qp}}(m\vect{k}) - E_{\mathrm{qp}}(n\vect{k})}{\varepsilon_{m\vect{k}} - \varepsilon_{n\vect{k}}} , \qquad m \neq n .
  \label{eq:E.3-vrescale}
\end{equation}
For a rigid shift \cref{eq:E.3-vrescale} is the scissor rescaling \cref{eq:3.4-vscis}. For a state-dependent correction it rescales each transition by its own factor. The velocity vertex enters the absorption of dark matter at small momentum transfer and the optical limit of $\rho_{mn}$ (\cref{sec:3.2}). \Cref{eq:E.3-vrescale} keeps the optical limit of the density vertex, $\vect{q}\cdot\vect{v}_{mn}/(E_m - E_n)$, equal to the optical limit in the uncorrected calculation, as it must be when the states are unchanged. The intraband velocities, the group velocities, follow from the dispersion of the corrected bands. Since the corrections are computed on a coarse grid, they are interpolated to fine grids with Wannier functions (\cref{app:E.5}), with the same gauge as every other interpolated quantity (\cref{sec:3.6}).

Effective masses change with the dispersion. The $\vect{k}\cdot\vect{p}$ argument of \cref{app:E.1} connects a mass at a direct gap to the gap, and the quasiparticle self-consistent calculation of GaAs shows the connection: it overestimates the gap, \SI{1.77}{eV} against \SI{1.52}{eV}, and the conduction-band mass at $\Gamma$, $0.077\,m_e$~\cite[Fig.~2 and text]{vanschilfgaarde2006quasiparticle} against the measured $0.067\,m_e$.

\subsection{Convergence and cost}

A $G_0W_0$ calculation is converged in the number of empty bands, the dielectric cutoff, the grid of crystal momenta and the frequency treatment, and these parameters are coupled. The number of empty bands $N_{\mathrm{b}}$ enters the sums for $R^0$ and for $\Sigma_{\mathrm{c}}$. The dielectric cutoff is the kinetic energy $E^{\varepsilon}_{\mathrm{cut}}$ below which the vectors $\vect{q} + \vect{G}$ of the matrix $\varepsilon_{\vect{G}\vect{G}'}$ are kept. The highest empty band included should have an energy of about the dielectric cutoff, since a plane wave $\vect{q} + \vect{G}$ resolves transitions into states up to that energy~\cite[Sect.~4]{deslippe2012berkeleygw}. In a free-electron count each plane wave with kinetic energy below $E$ contributes one band below $E$, so the number of bands below $E$ is about $N_{\mathrm{PW}}(E)$ of \cref{eq:E.1-npw}. For silicon, a dielectric cutoff of \SI{10}{Ry} keeps 144 vectors $\vect{G}$ and asks for about 150 bands; \SI{15}{Ry} keeps 265 and asks for about 270. The sum for the Coulomb-hole part of $\Sigma_{\mathrm{c}}$ converges much more slowly with $N_{\mathrm{b}}$ than the sum for $R^0$. A test that converges $R^0$ alone, or one parameter with the others held fixed, can therefore stop at a value far from the converged one~\cite[Sect.~4 and Fig.~8]{deslippe2012berkeleygw}. The grid of $\vect{k}$ and $\vect{q}$ points is the same for $R^0$ and for $\Sigma$. The divergence of $v_{\mathrm{C}}(\vect{q})$ at $\vect{q} \to 0$ is handled by averaging $v_{\mathrm{C}}$ and $W^{\mathrm{scr}}$ over the small cell around $\vect{q} = 0$ and by computing the head and wings of $\varepsilon_{\vect{G}\vect{G}'}$ (its $\vect{G} = \vect{G}' = 0$ element and the rest of the $\vect{G} = 0$ row and column) at $\vect{q} \to 0$ from the optical limit of the vertex. The frequency treatment is either a plasmon pole or a frequency grid. The bare exchange $\Sigma_{\mathrm{x}}$ is a separate case: it involves the unscreened interaction, whose $1/|\vect{q} + \vect{G}|^2$ falls off slowly, and its sum over $\vect{G}$ runs to the cutoff of the pair densities rather than to the dielectric cutoff.

The cost follows from the structure of the sums~\cite[Sect.~4]{deslippe2012berkeleygw}. The number of occupied and of empty bands, and the number of plane waves, each grow in proportion to the number of atoms $N_{\mathrm{at}}$. Computing the matrix elements \cref{eq:E.3-M} for all pairs of bands is $N_{\mathrm{at}}^2$ pairs times an FFT of size $N_{\mathrm{at}}\log N_{\mathrm{at}}$, so $N_{\mathrm{at}}^3\log N_{\mathrm{at}}$. Summing them into $R^0_{\vect{G}\vect{G}'}$ is $N_{\mathrm{at}}^2$ pairs for each of $N_{\mathrm{at}}^2$ pairs of $\vect{G}$ vectors, so $N_{\mathrm{at}}^4$, done as one dense matrix product; the inversion of $\varepsilon$ is $N_{\mathrm{at}}^3$ per $\vect{q}$ and per frequency. The self-energy of one state is $N_{\mathrm{at}}^3$, and all states in an energy window add a factor $N_{\mathrm{at}}$. A $G_0W_0$ calculation thus scales formally as the fourth power of the system size, against the third power of \cref{eq:E.1-cost}, with a prefactor set by the hundreds of empty bands. For the two-atom cell of silicon the cost remains modest. For the supercells of defects and for the dense grids of a finite-momentum exciton calculation, the $G_0W_0$ calculation is the step that limits what can be computed.

\begin{workedexample}{The silicon gap from three corrections}\label{we:E-si-gap}
A scissor, a hybrid functional and a $G_0W_0$ calculation each lead from the PBE calculation of \cref{we:E-si-bands} to a gap for a silicon threshold. The scissor to the measured gap at zero temperature uses $\Delta_{\mathrm{scis}} = \SI{0.568}{eV}$ at the PBE lattice constant, or \SI{0.613}{eV} at the measured one. It reproduces the measured gap by construction, includes the zero-point renormalization through the measurement, and is accurate to about \SI{0.1}{eV} for other bands~\cite{godby1988self}. The hybrid functional HSE gives \SI{1.22}{eV} at its own lattice constant (\cref{tab:E-functionals}). A one-shot $G_0W_0$ calculation from LDA gives \SI{1.10}{eV} converged, from a static lattice. Two corrections separate it from the measurement: spin--orbit coupling, which these calculations omit, and the zero-point renormalization. Spin--orbit coupling splits the sixfold valence-band maximum (with spin) into a fourfold level at $+\Delta_{\mathrm{SO}}/3$ and a twofold level at $-2\Delta_{\mathrm{SO}}/3$; the weighted mean, $4\times\Delta_{\mathrm{SO}}/3 - 2\times2\Delta_{\mathrm{SO}}/3 = 0$, is the level without spin--orbit coupling. The gap therefore closes by $\Delta_{\mathrm{SO}}/3$, \SI{15}{meV} in silicon and, with $\Delta_{\mathrm{SO}} = \SI{0.296}{eV}$, \SI{99}{meV} in germanium (\cref{sec:3.3}). The zero-point renormalization lowers the gap of a real crystal relative to the static lattice by about \SI{60}{meV}. For a dark-matter threshold in silicon, the three corrected gaps agree to about \SI{0.1}{eV}. The spread among them is the theory uncertainty of the threshold if the scissor is not anchored to a measurement, and \cref{we:3-mmin-gap} gives its effect on the minimum mass.
\end{workedexample}

The interaction of the excited electron with its hole, which a quasiparticle calculation omits, is treated by the Bethe--Salpeter equation of \cref{sec:4.9} and in \cref{app:E.6}.
